\documentclass[11pt]{../class/journalpaper}

\title{A Simple Example of a Journal Paper}

\author{
    Author A$^{1,*}$, Author B$^{1,2}$, Author C$^{2}$
    \\
    \small
    \parbox{0.95\linewidth}{
        \centering
        $^{1}$ \itshape Affiliation, University, Address, City, Postcode, State/Province, Country
        \\
        $^{2}$ \itshape Affiliation, University, Address, City, Postcode, State/Province, Country
        \\
        $^{*}$ Corresponding author: author@example.com
    }
}

\date{}

\begin{document}

\maketitle

\begin{abstract}
This document provides a complete example of the
\texttt{journalpaper} document class. It demonstrates every
feature defined by the class, including the title block,
abstract, keywords, section numbering, figure and table
captions, colored cross-references, and citation commands.
\end{abstract}


% cls 新建 keywords 环境：\textbf{Keywords: } 开头，\diyline 收束

\begin{keywords}
LaTeX; academic writing; journal paper; document class
\end{keywords}


% ============================================================
% 1. Introduction
% ============================================================

\section{Introduction}

This is a complete example demonstrating every feature of the
\texttt{journalpaper} document class.

The class preserves the standard structure of an academic
manuscript. Authors can use familiar commands such as
\verb|\title|, \verb|\author|, \verb|\date|, and
\verb|\maketitle|.

The class also provides predefined styles for sections,
figures, tables, equations, citations, and cross-references.


% ============================================================
% 2. Equation
% ============================================================
% cls 配置 cleveref 公式：
%   单个：Eq. (1)
%   多个：Eqs. (1) and (2)
%   区间：Eqs. (1)--(3)
% 全部通过 \cref / \Cref / \crefrange / \Crefrange 引用

\section{Equation}

Consider the following quadratic function:

\begin{equation}
    y = ax^2 + bx + c.
    \label{eq:quadratic}
\end{equation}

For $a=1$, $b=2$, and $c=1$, \cref{eq:quadratic} becomes

\begin{equation}
    y = x^2 + 2x + 1.
    \label{eq:specific}
\end{equation}

A third equation is added for range and multiple tests:

\begin{equation}
    y = (x+1)^2.
    \label{eq:third}
\end{equation}

% --- 单个公式引用 ---
Equation \cref{eq:quadratic} is the general form, and
\cref{eq:specific} is a special case.

% --- \cref 单个公式 ---
The class also supports \cref{eq:third}.

% --- \Cref 单个公式（句首） ---
\Cref{eq:quadratic} is the most general form.

% --- 多个公式引用（离散） ---
See \cref{eq:quadratic,eq:specific} for the general and
specific cases.

% --- 连续区间公式引用 ---
Equations \crefrange{eq:quadratic}{eq:third} cover the full
range.

% --- \Crefrange 句首区间 ---
\Crefrange{eq:quadratic}{eq:third} cover the full range.


% ============================================================
% 3. Table
% ============================================================
% cls 配置 cleveref 表格：
%   单个：Table 1
%   多个：Tables 1 and 2
%   区间：Tables 1--3
% 全部通过 \cref / \Cref / \crefrange / \Crefrange 引用

\section{Table}

\cref{tab:example} gives several example values calculated
from \eqref{eq:specific}.

\begin{table}[htbp]
    \centering

    \caption{Example values of the quadratic function.}
    \label{tab:example}

    \begin{tabular}{ccccc}
        \toprule
        $a$ & $b$ & $c$ & $x$ & $y$ \\
        \midrule
        1.0 & 2.0 & 1.0 & $-2$ & 1 \\
        1.0 & 2.0 & 1.0 & $-1$ & 0 \\
        1.0 & 2.0 & 1.0 & $0$  & 1 \\
        1.0 & 2.0 & 1.0 & $1$  & 4 \\
        1.0 & 2.0 & 1.0 & $2$  & 9 \\
        \bottomrule
    \end{tabular}
\end{table}

A second table is added for multiple and range tests.

\begin{table}[htbp]
    \centering

    \caption{A second example table.}
    \label{tab:second}

    \begin{tabular}{cc}
        \toprule
        Item & Value \\
        \midrule
        A & 1 \\
        B & 2 \\
        \bottomrule
    \end{tabular}
\end{table}

A third table is added for range tests.

\begin{table}[htbp]
    \centering

    \caption{A third example table.}
    \label{tab:third}

    \begin{tabular}{cc}
        \toprule
        Item & Value \\
        \midrule
        C & 3 \\
        D & 4 \\
        \bottomrule
    \end{tabular}
\end{table}

% --- 单个表格引用 ---
\Cref{tab:example} lists the values.

% --- 多个表格引用（离散） ---
\cref{tab:example,tab:second} give complementary data.

% --- 连续区间表格引用 ---
\crefrange{tab:example}{tab:third} cover all tables.

% --- \Crefrange 句首区间 ---
\Crefrange{tab:example}{tab:third} cover all tables.


% ============================================================
% 4. Figure
% ============================================================
% cls 配置 cleveref 图：
%   单个：Figure 1
%   多个：Figures 1 and 2
%   区间：Figures 1--3
% 全部通过 \cref / \Cref / \crefrange / \Crefrange 引用

\section{Figure}

\cref{fig:example} demonstrates the standard
\texttt{figure} environment.

\begin{figure}[htbp]
    \centering

    \fbox{%
        \rule{0pt}{45mm}%
        \rule{0.8\linewidth}{0pt}%
    }

    \caption{A placeholder figure.}
    \label{fig:example}
\end{figure}

A second figure is added for multiple and range tests.

\begin{figure}[htbp]
    \centering

    \fbox{%
        \rule{0pt}{45mm}%
        \rule{0.6\linewidth}{0pt}%
    }

    \caption{A second placeholder figure.}
    \label{fig:second}
\end{figure}

A third figure is added for range tests.

\begin{figure}[htbp]
    \centering

    \fbox{%
        \rule{0pt}{45mm}%
        \rule{0.4\linewidth}{0pt}%
    }

    \caption{A third placeholder figure.}
    \label{fig:third}
\end{figure}

% --- 单个图引用 ---
\Cref{fig:example} shows a placeholder.

% --- 多个图引用（离散） ---
\cref{fig:example,fig:second} show two placeholders.

% --- 连续区间图引用 ---
\crefrange{fig:example}{fig:third} show three placeholders.

% --- \Crefrange 句首区间 ---
\Crefrange{fig:example}{fig:third} show three placeholders.


% ============================================================
% 5. Section Formatting
% ============================================================
% cls 重定义章节编号：
%   section:       1.
%   subsection:    1.1.
%   subsubsection: 1.1.1.

\section{Section Formatting}

The section numbering is automatically formatted by the class.

\subsection{Subsection}

This is a subsection. Its number appears in the form
``5.1.''.

\subsubsection{Subsubsection}

This is a subsubsection. Its number appears in the form
``5.1.1.''.


% ============================================================
% 6. Citation
% ============================================================
% cls 加载 natbib，参考文献样式由 \bibliographystyle 指定
% 引用命令：\citet / \citep

\section{Citation}

\LaTeX{} is widely used for scientific and technical document
preparation \citep{lamport1994}.


% ============================================================
% 7. Conclusion
% ============================================================

\section{Conclusion}

This example demonstrates every feature of the
\texttt{journalpaper} document class: the title block,
abstract, keywords, section numbering, figure and table
captions, colored cross-references for single, multiple, and
ranged figures, tables, and equations, as well as citations.


% ============================================================
% References
% ============================================================

\bibliographystyle{apalike}
\bibliography{references}

\end{document}